\documentclass[11pt]{amsart}

\usepackage{amsmath, amssymb, amsthm}
\usepackage{mathtools}
\usepackage{thmtools}
\usepackage{enumitem}
\usepackage{tikz-cd}
\usepackage{hyperref}
\usepackage[nameinlink]{cleveref}

% % Theorem environments
\declaretheorem[numberwithin=section]{theorem}
\declaretheorem[style=plain, sibling=theorem]{lemma, proposition, corollary}
\declaretheorem[style=definition]{definition, example, condition}
\declaretheorem[style=remark]{remark, note, observation, todo}


\title{April 13, 2026}


\begin{document}
\maketitle

We first configure two skills of Codex.

\begin{enumerate}
    \item \texttt{writing-feedback}: This skill read the raw log file, and provide feedback on it.
    The skill provide grammatical or mathematical feedback depending on the type of the given log code.
    \item \texttt{arXiv-review-routine}: This skill fetchs arxiv's updated papers and filter papers related to author's interests.
    After fetching the papers, the user must copy the file generated in \texttt{derived/arxiv/review/generated} with a file name of today, and paste it in \texttt{derived/arxiv/review/checked}.
    After that, the user must read the file and check whether each paper is interested or not.
    This information will be used for future arXiv paper filter criteria.
\end{enumerate}

After configuring the above skills, I recognized to implement automation of those skills.

\bibliographystyle{amsalpha}
% \nocite{*}
\bibliography{bibliography}
\end{document}
